\section{Detect Cycle in an Undirected Graph}
\label{sec:graph-cycle-undirected}

\problemstatement{Given an undirected graph, determine whether it
contains a \textbf{cycle}.}

\begin{patternbox}
\emph{``Does a cycle exist''} is a traversal that carries one extra
piece of state along: the \textbf{parent} node (the node the current
node was reached from). Reaching an already-visited neighbor that is
\textbf{not} the parent means there are two distinct paths to that
neighbor -- which is exactly what a cycle is.
\end{patternbox}

\subsection*{Why checking against the parent specifically is necessary}

In an \emph{undirected} graph, every edge $(u,v)$ is stored twice: once
as $u \to v$ and once as $v \to u$. So when DFS or BFS moves from $u$
to $v$, $v$'s neighbor list contains $u$ again -- without excluding the
parent, this single edge would be misreported as a cycle every time.
Excluding only the \emph{immediate} parent (not all previously visited
nodes in general) is correct because any \emph{other} already-visited
neighbor reached again is a genuinely different path back to that
node, while the parent connection is simply the edge just traversed in
reverse.

\subsection*{Approach 1: DFS}

\begin{lstlisting}
#include <bits/stdc++.h>
using namespace std;

bool dfsHasCycle(int node, int parent, vector<vector<int>>& adj, vector<bool>& visited) {
    visited[node] = true;

    for (int neighbor : adj[node]) {
        if (!visited[neighbor]) {
            if (dfsHasCycle(neighbor, node, adj, visited)) return true;
        } else if (neighbor != parent) {
            return true; // visited, and not the edge we just came from
        }
    }
    return false;
}

bool hasCycleDFS(int n, vector<vector<int>>& adj) {
    vector<bool> visited(n, false);
    for (int i = 0; i < n; i++) {
        if (!visited[i] && dfsHasCycle(i, -1, adj, visited)) {
            return true;
        }
    }
    return false;
}

int main() {
    int n = 4;
    vector<vector<int>> adj(n);
    auto addEdge = [&](int u, int v) { adj[u].push_back(v); adj[v].push_back(u); };
    addEdge(0, 1); addEdge(1, 2); addEdge(2, 3); addEdge(3, 0);

    cout << (hasCycleDFS(n, adj) ? "true" : "false") << endl; // true
    return 0;
}
\end{lstlisting}

\begin{complexitybox}[Approach 1 Complexity]
\textbf{Time.} $O(V+E)$.
\textbf{Space.} $O(V)$ for \textit{visited} and the recursion stack.
\end{complexitybox}

\subsection*{Approach 2: BFS}

The same parent-tracking idea, carried in the queue as a
$(\textit{node}, \textit{parent})$ pair instead of as a recursion
argument.

\begin{lstlisting}
#include <bits/stdc++.h>
using namespace std;

bool bfsHasCycle(int src, vector<vector<int>>& adj, vector<bool>& visited) {
    queue<pair<int,int>> q; // {node, parent}
    visited[src] = true;
    q.push({src, -1});

    while (!q.empty()) {
        auto [node, parent] = q.front();
        q.pop();

        for (int neighbor : adj[node]) {
            if (!visited[neighbor]) {
                visited[neighbor] = true;
                q.push({neighbor, node});
            } else if (neighbor != parent) {
                return true;
            }
        }
    }
    return false;
}

bool hasCycleBFS(int n, vector<vector<int>>& adj) {
    vector<bool> visited(n, false);
    for (int i = 0; i < n; i++) {
        if (!visited[i] && bfsHasCycle(i, adj, visited)) {
            return true;
        }
    }
    return false;
}

int main() {
    int n = 4;
    vector<vector<int>> adj(n);
    auto addEdge = [&](int u, int v) { adj[u].push_back(v); adj[v].push_back(u); };
    addEdge(0, 1); addEdge(1, 2); addEdge(2, 3); addEdge(3, 0);

    cout << (hasCycleBFS(n, adj) ? "true" : "false") << endl; // true
    return 0;
}
\end{lstlisting}

\subsection*{Dry run}

\dryrun{Take the $4$-cycle $0{-}1{-}2{-}3{-}0$, DFS starting at $0$.}

$\textit{dfs}(0,-1)$: visit $0$. Neighbor $1$: unvisited, recurse
$\textit{dfs}(1,0)$: visit $1$. Neighbor $0$: visited, but
$0=\textit{parent}$, skip. Neighbor $2$: unvisited, recurse
$\textit{dfs}(2,1)$: visit $2$. Neighbor $1$: visited, $=\textit{parent}$,
skip. Neighbor $3$: unvisited, recurse $\textit{dfs}(3,2)$: visit $3$.
Neighbor $2$: visited, $=\textit{parent}$, skip. Neighbor $0$: visited,
and $0 \neq \textit{parent}(=2)$ -- \textbf{cycle detected}, return
\texttt{true} all the way up.

\begin{complexitybox}[Approach 2 Complexity]
\textbf{Time.} $O(V+E)$.
\textbf{Space.} $O(V)$ for \textit{visited} and the queue.
\end{complexitybox}

\begin{takeawaybox}
Both traversals solve this with the identical idea -- carry the parent
along, flag any visited non-parent neighbor as a cycle -- differing
only in whether that extra state rides on the recursion stack (DFS) or
in the queue alongside each node (BFS). This same parent-carrying
trick reappears, adapted, in the Topological Sort chapter's directed-graph
cycle detection.
\end{takeawaybox}
